\section{An exact renormalization map}
\label{rg:section}

The use of local backgrounds and small-field integration has close
constructive predecessors
\cite{BalabanFlow,BalabanVariational,DybalskiStottmeisterTanimotoVariational}.
The exact observation kernel and the separate treatment of rough
configurations in the retained-density class are specified below; the
cited constructions do not supply closure or terminal-coupling comparison.

The estimates in this section concern the original lattice measure at
finite cutoff.  The renormalization map will be an exact conditional
integration.  Its purpose is to compare two cutoffs after they have reached
the same value of the running coupling.  The terms not described by that
coupling must therefore be controlled as functions, including their
dependence on the cutoff.

Section~\ref{rg:class} defines the retained class and states closure.
Section~\ref{rg:gaussian} constructs the positive energy reserve and
the Gaussian integral. Sections~\ref{rg:prepared} and~\ref{rg:connected}
bound joint products and reassemble the exact density; the explicit
factor ownership and support tables are in
Sections~\ref{app:factor-ownership} and~\ref{app:read-sets}, and the
product estimate is proved in Appendix~\ref{supp:joint-majorants}.
Sections~\ref{rg:canonical} and~\ref{rg:error} then contract the
canonical coefficients and regular error, using the fixed norm of
Appendix~\ref{supp:canonical-fixed-norm}. Section~\ref{rg:normalization}
restores the normalization and closes the induction.

We identify $S^3$ with the unit quaternions.  Multiplication on either side
is an isometry, and the stabilizer of the identity acts by all rotations
on the imaginary quaternions $\mathbb R^3$.  Haar measure is normalized to
have total mass one.  We use the averaging operator $Q$, the precisions
$K_h$, and the operators
\[
 \mathsf A_h=\binom{\sqrt{c_0}\,\mathrm{Id}}{\ell_h},\qquad
 P_h,\quad C_h,\quad \mathcal M_h,\quad\mathcal N_h
\]
of Lemmas~\ref{free:bounds} and~\ref{lem:free-stack}.  In particular,
$K_h=d^*\mathsf A_h^*\mathsf A_h d$, and the constants in the weighted
first and second fine-difference bounds for $P_h$ are independent of
sufficiently large $L$.  Constants denoted by $C_L,c_L>0$ may depend on
the now fixed blocking factor.  Constants used to choose $L$ will be
explicitly distinguished from these.

\subsection{Scales, observation, and the class of densities}
\label{rg:class}

Choose a dyadic integer $L$, later sufficiently large, and put
\begin{equation}
 \begin{aligned}
 \gamma&=\frac{\log L}{\pi},&
 H_j&=H+\gamma j,&
 \mathfrak m_j&=\left\lceil(\log H_j)^2\right\rceil,\\
 p_j&=(\log H_j)^{P_0},&
 t_j&=H_j^{-1/2}p_j,&
 T_j&=H_j^{-.49}.
 \end{aligned}
 \label{rg:scales}
\end{equation}
The exponent $P_0$ and then $H$ will be fixed below.  A run of depth $N$
uses layers $j=N,N-1,\ldots,0$ and free history $h=N-j$.  At a step its
precise coupling $b$ is required to lie in $[H_j/2,2H_j]$.  We write
$g=b^{-1/2}$ and omit $j,h$ temporarily.  Within a single step,
$M=\mathfrak m_j$ and $M'=\mathfrak m_{j-1}$.  Thus $t\asymp gp$.  All
geometric thresholds in a comparison use $H_j$, not the varying $b$.

On a square torus of side $mL^j$, take $m$ dyadic and
$m\ge C_{\rm adm}\mathfrak m_0$.  The constant $C_{\rm adm}$ is fixed sufficiently
large for the finite neighborhoods below.  Increasing $H$ gives
$\mathfrak m_j/\mathfrak m_{j-1}\le2$; hence these period conditions hold simultaneously
at every layer.  Bulk formulas mean the same formulas with their
absolutely summable position lists on $\mathbb Z^2$.

The observation from $q$ to one new spin $V_Y$ in each block is
\begin{equation}
 \begin{aligned}
 \mathcal K_b(dV\mid q)&=
 \mathbb E_{\tau}\!\left[
 \prod_Y\frac{\exp\{-ba|V_Y-\widehat m_Y(q,\tau_Y)|^2/2\}}
 {z(ba)}\,dV_Y\right],\\
 z(ba)&=\int_{S^3}e^{-ba|V-1|^2/2}\,dV,
 \qquad m_Y=Qq(Y).
 \end{aligned}
 \label{rg:observation}
\end{equation}
Here $a$ is the fixed positive observation precision in the free
construction.  If $|m_Y|\ge1/2$, set $\widehat m_Y=m_Y/|m_Y|$.
Otherwise the independent tag $\tau_Y$ samples a constituent spin
with the weights of $Q$, and $\widehat m_Y(q,\tau_Y)$ is that spin.
The expectation in \eqref{rg:observation} averages these tags.
They may be included at every block, even where they are unused.  The denominator
in \eqref{rg:observation} is independent of the unit vector
$\widehat m_Y$, so this is a probability kernel.

For an edge $e$, write $r_e=|d_eq|$.  Fix $0<d_0<c_0/16$ and a constant
$R_*$ sufficiently large compared with the exponential localization
length of $\ell_h$.  Define
\begin{align}
 S_t(r)&=\begin{cases}c_0r^2/2,&r\le t,\\d_0tr,&r>t,\end{cases}
 \nonumber\\
 m_e(q)&=\mathbf1\{r_f\le t e^{\operatorname{dist}(e,f)/R_*}
                                  \text{ for all }f\},\nonumber\\
 \mathcal E_{h,t}(q)&=\sum_eS_t(r_e)
       +\frac12\sum_em_e(q)| (\ell_hdq)_e|^2.
 \label{rg:kinetic}
\end{align}
The mask reads only edges within $O(\log(1/t))$ of $e$, since chord
lengths are at most two.  Moreover $m_e=1$ implies
$| (\ell_hdq)_e|\le Ct$.  The masks and $S_t$ may be discontinuous;
they will never be differentiated as functions of a moving output
configuration.

We next specify the locality norms.  Each label carries its complete
support and a connected record of cubes of radius $\mathfrak m_j$ covering that
support, together with the paths used to join them.  Its load is an
integer $s\ge1$.  The conventions are chosen so that the support has at
most $C\mathfrak m_j^2s$ sites and its record has length at most $C\mathfrak m_js$.
Every argument, mask, eligibility test, negative status query, and
window determining a formula is included.  On a torus the record retains
the lifted displacements before folding.  Intersecting records can be
joined with load at most a fixed multiple of their total load.

Projection of a record to the next lattice, followed by any of the
fixed neighborhoods used below, has a record satisfying
\begin{equation}
 4s/L\le s'\le D_*(1+s/L).
 \label{rg:projection}
\end{equation}
The upper bound follows by breaking its walk into pieces of length
comparable with $L\mathfrak m_{j-1}$; the lower bound is imposed by retaining
unused load if necessary.  A record of load less than
$c_*mL^j/\mathfrak m_j$ is called short.  Choose $c_*$ small enough that every
short complete record has an injective lift.  Its formula must coincide
with that of the corresponding bulk label.  All other records are
called winding records, whether or not their displayed support itself
wraps.  This convention also applies to a constant or a short function
whose derivation used a long record.

For a connected graph $X$ let
$1_X(q)=\mathbf1\{r_e\le t_j\text{ for all }e\in X\}$.
At an anchor $o$, use relative coordinates
\[
 y_o(x)=\operatorname{Im}(q_xq_o^{-1}),\qquad
 S_o(q)=\frac14\sum_{|e|=1}|y_o(o+e)|^2.
\]
Fix a basis of each space of invariant tensors
$((\mathbb R^3)^{\otimes n})^{SO(3)}$, $2\le n\le6$.
In particular this includes the alternating cubic tensor.  A canonical
polynomial is a homogeneous tensor polynomial in the $y_o(x)$, with an
anchored joining path of length $u$ comparable with
$1+\sum_x|x-o|$, counted with multiplicity.  Its coefficient norm is
\begin{equation}
 \|A_n\|_\sigma=
 \sup_o\sum_{\text{labels at }o}e^{\sigma u}(1+u)^2
                                  |a_{\text{label}}|.
 \label{rg:canonical-norm}
\end{equation}
Take $\sigma>0$ sufficiently small for the finitely many uses of the
free exponential bounds; it is fixed before $L$.  The seven canonical
slots, in their order, are
\begin{equation}
 bP_4,\quad I_2,\quad bP_5,\quad I_3,\quad
 bP_6,\quad I_4,\quad b^{-1}J_2.
 \label{rg:slots}
\end{equation}
The subscript records homogeneous field degree, and
$\mathbf P=(P_4,I_2,P_5,I_3,P_6,I_4,J_2)$ denotes the seven coefficient
lists without their displayed powers of $b$. Their norms have fixed caps
$B_1,\ldots,B_7$, chosen successively below.

Each quadratic label is grouped into its dihedral orbit at the anchor
and compensated by a multiple of $S_o$, so that its Hessian vanishes on
unit affine fields at the identity.  This means its second derivative
on $q_x=\exp\{\theta v\,\widehat e\cdot(x-o)\}$ is zero, for a unit
imaginary quaternion $v$ and a coordinate unit vector $\widehat e$.
Spin and spatial symmetry then imply vanishing on every pair of affine
tangent fields.  Compensation coefficients are included in
\eqref{rg:canonical-norm}.  All polynomials are formed before folding.
For $u\le \mathfrak m_j$ their mask graph may be a ball of radius
$C_\chi \mathfrak m_j$ at $o$.  Longer paths are padded by fixed multiples of
$\mathfrak m_j$.  Their load is at most $C(1+u/\mathfrak m_j)$.

A regular error $f_X$ is defined and $C^k$, with $k=8$, on
$\{r_e<4t_j:e\in X\}$.  Its norm $[f_X]_{k,j}$ is the sum of the sup
norms of its derivatives through order $k$, using independently in
each derivative slot the bounds
\begin{equation}
 |v_x|\le t_j(1+\operatorname{dist}(x,o))^8.
 \label{rg:directions}
\end{equation}
Derivatives are taken in
$q_x(\theta)=\exp(\sum_i\theta_iv_{i,x})q_x$.
On a torus use physical distances.  The error norm and cap are
\begin{equation}
 \|f\|_{k,j,A}:=\sup_o\sum_{X\text{ at }o}e^{As_X}[f_X]_{k,j}
       \le\delta_j:=H_j^{-2.05},\qquad A=64.
 \label{rg:error-norm}
\end{equation}
Each bulk or nonwinding error has zero value and zero affine Hessian
at the identity.  Its first derivative is zero by conjugation
invariance.  Winding errors need not have these normalizations.

Finally a covering label $\ell$ has a bounded measurable weight
$k_\ell(q)$, complete support $P_\ell$, and a nonempty inventory
$D_\ell$ of bonds in that support.  Its weight vanishes unless
$D_\ell\subset D(q):=\{e:r_e>t_j\}$.  Set
\begin{equation}
 \Xi(q)=\sum_{\substack{\Gamma\text{ has disjoint supports}\\
                  \bigsqcup_{\ell\in\Gamma}D_\ell=D(q)}}
                    \prod_{\ell\in\Gamma}k_\ell(q),
 \qquad
 \|k\|_{j,A}:=\sup_x\sum_{\ell:x\in P_\ell}
                  e^{As_\ell}\|k_\ell\|_\infty
       \le w_j:=e^{-p_j^{1/4}}.
 \label{rg:covering-gas}
\end{equation}
In particular $\Xi=1$ if $D(q)$ is empty.  This is a mandatory cover,
not an optional polymer gas.  No sign or differentiability assumption
is made on the individual weights.

All these lists are translation covariant and have the spin and
spatial symmetries just specified.  The class consists of the densities
\begin{equation}
 e^{v|\Lambda_j|}
 \exp\left\{-b\mathcal E_{h,t_j}
       -\sum_{\text{slots}}1_X P_X-\sum_X1_Xf_X\right\}\Xi.
 \label{rg:density}
\end{equation}
The scalar $v$ is the same bulk scalar on every admitted period.
Assignments of anchors are made equivariantly: distribute a term equally
among its permissible anchors, retaining identical complete supports
on the resulting copies.  Splitting an additive term this way preserves
its sum; splitting a covering weight preserves its gas, since copies
with the same nonempty support are mutually incompatible.  Average
additive anchored coefficient and error lists over the finite dihedral
group.  Their full sums are invariant, so this also preserves the density.
If different orbit terms initially have different masks, first strengthen
them to their common orbit graph and retain the exact covering correction,
as in Section~\ref{rg:connected}.  Eligibility and internal compatibility
tests are carried with each label, not averaged independently.  All
geometric rules and free kernels commute with these symmetries.
Finally, after removing the total affine coefficient, orbitwise quadratic
compensation subtracts multiples of $S_o$ whose sum is zero at every
anchor; it changes only the representation.  These conventions justify
both the individual orbit normalization and the invariant Hessian test.
Masked expressions are defined to be zero off their mask; no value
of an undefined extension is used there.  The estimates below also
apply to signed densities of this form.  Positivity is used later only
for the actual probability measures obtained from
\eqref{rg:observation}.

\begin{theorem}[Exact step and its comparison estimates]
\label{rg:closure}\label{thm:rg-closure}
The constants can be chosen in the order described below so that
integration by \eqref{rg:observation}, from an input with
$b\in[H_j/2,2H_j]$, gives an exact representation
\eqref{rg:density} on layer $j-1$, with free history $h+1$, all the
renewed norm bounds, and
\begin{equation}
 b'=b+\alpha_h(\mathbf P)+\frac{\kappa_h(\mathbf P)}b+\Delta,
 \qquad |\alpha_h|+|\kappa_h|\le C_L,
 \qquad |\Delta|\le C_LH_j^{-1.05}.
 \label{rg:coupling-step}
\end{equation}
The canonical map and $\alpha_h,\kappa_h$ depend only on the input
canonical vector $\mathbf P$ and the free history.  They do not depend
on the precise coupling, the cutoffs, the errors, the covering gas,
or the period.  The output is iterable if its coupling belongs to the
next admitted band; this additional fact is proved for the trajectories
used below.

For two inputs on common label lists and the same reference layer,
let $\lambda=b_2-b_1$.  Let $u$ be a fixed positively weighted sum of
the seven coefficient discrepancies, the error discrepancy
$\delta_j^{-1}\|f_2-f_1\|_{k-1,j,A}$, and the gas discrepancy
$w_j^{-1}\|k_2-k_1\|_{j,A}$.  The weights are independent of $j$ and
the periods.  The outputs have common compatible lists and satisfy
\begin{align}
 u'&\le q u+C_LH_j^C\epsilon_h+C_LH_j^{-b_*}|\lambda|,
 \label{rg:shape-comparison}\\
 |\lambda'-\lambda|
    &\le C_Lu+C_LH_j^C\epsilon_h+C_LH_j^{-b_*}|\lambda|,
 \label{rg:coupling-comparison}
\end{align}
where $L^{-1}<q<1$, $b_*>0$, and
$\epsilon_h=L^{-\min(h_1,h_2)}$ for unequal free histories
(zero for equal histories).  These estimates include bulk and all
admitted periods.
\end{theorem}

To prove closure, we first isolate regions with large gradients and
retain a positive energy cost for them. Gaussian completion on the
complement produces local fluctuation factors. We then sum connected
factors, extract the seven canonical coefficients, and return the
remaining terms to the error and covering classes.

Appendix~\ref{app:rg-geometry} fixes the profiles, row ownership and
complete read sets used below. It supplies an explicit core factor and
proves the derivative-support and factorization assertions.

\subsection{The energy identity and the massive Gaussian}
\label{rg:gaussian}

Truncate the primitive kernels at radius
$r_*=\lfloor M/10^4\rfloor$ in their respective localization units,
reducing this by a fixed factor if necessary for the finite compositions
below.  Use a subscript $s$ for these numerical kernels.  The omitted
paths have exponentially small weighted sums, including their free
history differences.  Retain those paths as separate factors.

Use stack variables
\[
 X=(\sqrt{c_0}\,dq,\ell_sdq),\qquad
 Y=\sqrt a(V-\widehat m),\qquad
 X'=(\sqrt{c_0}\,d'V,\ell'_sd'V),
\]
and put
\[
 \zeta=(\sqrt{c_0}\,\operatorname{clip}_{t'}d'V,
                                     m'\ell'_sd'V),
 \qquad(\nu,w)=\mathcal M_s\zeta.
\]
Here clipping changes a vector of length greater than $t'$ to the
vector of length $t'$ in the same direction.  The free estimates and
the masks give, entrywise,
\begin{equation}
 |\zeta|+|w|\le Ct,\qquad |\nu|\le Ct/L.
 \label{rg:stack-sizes}
\end{equation}
The numerical input kinetic energy plus observation energy minus
the numerical output kinetic energy is the sum of the following rows:
\begin{align}
 \mathcal F_e&=S_t(r_e)+\tfrac12m_e|X_{2,e}|^2
                      -\nu_e\cdot X_e+\tfrac12|\nu_e|^2,
 &\mathcal B_Y&=\tfrac12|Y_Y-w_Y|^2,\nonumber\\
 \mathcal C_E&=\zeta_E\cdot X'_E-\tfrac12|\zeta_E|^2
                  -S_{t'}(|d'_EV|)-\tfrac12m'_E|X'_{2,E}|^2,
 \nonumber\\
 \mathcal U_E&=\zeta_E\cdot[\mathcal M_s^*(X,Y)-X']_E,
 &\mathcal N_i&=\tfrac12| (\mathcal N_s\zeta)_i|^2,\nonumber\\
 \mathcal Z_E&=\tfrac12\zeta_E\cdot
     [(\mathrm{Id}-\mathcal M_s^*\mathcal M_s-
                       \mathcal N_s^*\mathcal N_s)\zeta]_E.
 \label{rg:ledger}
\end{align}
This is an algebraic identity: the cross terms cancel because
$(\nu,w)=\mathcal M_s\zeta$, and the displayed defect cancels the
remaining quadratic terms in $\zeta$.  The free stack identities give
\begin{equation}
 |\mathcal Z_E|\le C_Le^{-c_LM}t^2,
 \quad
 |\mathcal U_E|\le C_Lte^{-c_LM}
       +Ct\sum_Y e^{-c\operatorname{dist}(E,Y)}|m_Y-\widehat m_Y|.
 \label{rg:ledger-errors}
\end{equation}

Insert a smooth decomposition $G_e+(1-G_e)$ at every input edge,
where $G_e=1$ for $r_e\le t/4$ and $G_e=0$ for $r_e\ge t/2$.
At every observation block use an analogous decomposition for
$|V_Y-\widehat m_Y|$, with thresholds $Wt$ and $2Wt$; $W$ is a
fixed sufficiently large constant.  A complement is a primary seed.
Every true output bad edge is also a mandatory primary seed.  Join
seeds whose radius-$50M$ footprints meet.  For each component freeze
the spins within distance $3M$ of its seeds and assign to it the
numerical rows within distance $5M$.  Frozen spins retain their Haar
integration.  The remaining spins have the principal coordinates
\[
 T_x=V_{[x]},\qquad \xi_x=\log(q_xT_x^{-1}),\qquad
 q_x=\exp(\xi_x)T_x,
 \quad \xi_x\in\operatorname{Im}\mathbb H.
\]
The retained sector tests imply $|\xi_x|\le C_Lt<\pi/4$.
Thus this logarithm is unique; an arbitrary preimage under the
exponential map is not allowed.

Select all $\mathcal F,\mathcal B$ rows farther than $2.5M$ from
every seed.  On the full sector their masks are enabled, clipping is
absent, their mean is on its smooth branch, and all their spins are
within $C_Lt$ of their block reference.  Rotate a row to a nearby
output reference and let
\[
 D_0=(\mathsf A_s d,-\sqrt aQ),\qquad
 F_i^0(V)=\operatorname{Im}(R_i(0,V)V_{o(i)}^{-1}),
 \quad R=(X-\nu,Y-w).
\]
Use the affine row $D_0\xi+F^0$.  A frozen angle is smoothly cut off
outside a slightly larger $C_Lt$ chart.  For an output-inventory-free
call, its $V$ extension fixes every queried output status to no flag,
replaces clipping by the identity and every good output mask by one,
and extends that analytic formula through the required graph gradients
$<4t'$.  Cut it off smoothly only on a larger graph domain.  Equality
with the actual expression is recovered on the outgoing regular mask.
No extension of a discontinuous output predicate through its jump is
intended.  Calls with output inventory instead retain their fixed output
statuses and are not field-differentiated.  These conventions give the
bounded extensions used below on arbitrary frozen spins and output data.
In exterior variables the selected affine rows are $D\xi+F$.

A \emph{protector} is a smooth cutoff in an exterior coordinate $\xi_x$,
equal to one for $|\xi_x|\le\pi/2$ and supported strictly inside
$|\xi_x|<\pi$. It remains attached to every factor reading that protected
coordinate. The attachment rules are given in Section~\ref{rg:prepared}.

\begin{lemma}[Core reserve]
\label{rg:core-reserve}
On the retained sector tests around a component with $s_c$ seeds,
using these principal coordinates,
\begin{equation}
 \sum_{\rm assigned}(\mathcal F+\mathcal B+\mathcal C+
                           \mathcal U+\mathcal N+\mathcal Z)
 -\frac12\sum_{\rm selected,assigned}|D\xi+F|^2
                   \ge c t^2s_c.
 \label{rg:reserve}
\end{equation}
The same estimate holds on the no-output-flag extension of a
component with no output seeds.  After extending the Gaussian
coordinates beyond the original chart, it holds on the attached
protector supports, including the supports of their
fixed-order derivatives and of the retained sector-profile
derivatives.  No estimate for an arbitrary exponential preimage
outside those supports is asserted.
\end{lemma}
\begin{proof}
On a good input edge the direct part of $\mathcal F$ is a square.
At an input seed $r_e\ge t/4$ it pays $ct^2$ by
\eqref{rg:stack-sizes}, after choosing $L$ large.  If $r_e>t$ it
pays $ctr_e$.  The second-stack part is nonnegative unless its mask
fails.  A failed mask has a witness $f$ with $r_f>t$ in this same
component.  For each witness there are at most
$C(1+\log(r_f/t))^2\le Cr_f/t$ failed rows.  Their contributions from
chords at most $t$ cost at most $Ct\sum_fr_f/L$; contributions from
larger chords satisfy the same bound by the column sum of $\ell_s$.
Absorb them into the direct reserves.  Noise seeds pay $ct^2$ when
$W$ is large.  On an output flag the direct part of $\mathcal C$ is
\[
 (c_0-d_0)t'r-\tfrac12c_0t'^2
                    \ge(c_0/2-d_0)t'r,
\]
and its second-stack part cancels exactly.  The row
$\mathcal N$ is nonnegative.

For the remaining mean error, coordinate paths in a block give
\[
 |m_Y-\widehat m_Y|\le \frac CL\sum_{e\subset Y}r_e.
\]
For the fallback branch, $|m_Y|<1/2$ already forces the right side
to be bounded below by a positive constant, so this estimate holds
for every sampled constituent spin.  If every block chord is at
most $T$, the smooth branch is used and the stronger bound is
$C_LT^2$.  In a block containing a chord greater than $T$, charge
the large-chord terms in \eqref{rg:ledger-errors} to the same direct
reserves with their $L^{-1}$ gain.  Its remaining $C_Lt^2$ cost is
paid by the stronger $ctT$ witness.  All other mean and numerical
losses total at most
$C_LM^Ds_c(tT^2+te^{-c_LM})$.

Selected assigned rows were used only for nonnegativity in these
payments.  On their retained defaults all angles are the genuine
small angles, and the actual row differs from its affine row by
$C_LM^Dt^2$.  Subtracting its affine square therefore costs at most
$C_LM^Dt^3$ per row.  These losses are $o(t^2s_c)$, since
$T^2/t=H_j^{-.48}/p$, $tM^D\to0$, and
$M^De^{-c_LM}/t\to0$.

All witnesses and defaults used here lie within the retained $14M$
neighborhood of the seeds.  A derivative of a complement still has
$r_e\ge t/4$ or noise at least $Wt$; a derivative of a default
retains its upper support bound.  The attached protectors keep every
exterior argument in the injectivity ball.  Together with the
retained physical-spin tests this selects precisely its small
principal logarithm, also on all the stated derivative supports.
Thus the same argument applies there.  In a component with no output seeds use
the unclipped, positive-mask output formula through $4t'$; only the
fixed constants in the preceding estimates change.
\end{proof}

The exterior quadratic $A=D^*D$ is the principal restriction of
\[
 d^*\mathsf A_s^*\mathsf A_s d+aQ^*Q.
\]
Indeed every row incident on an exterior column is selected.
For a vector extended by zero into the frozen region, block Poincare
gives
\[
 \|u\|_2^2\le CL^2(\|du\|_2^2+\|Qu\|_2^2).
\]
Hence $c_L\mathrm{Id}\le A\le C_L\mathrm{Id}$, uniformly in the
primary geometry and the period.

For an exterior site $x$ let $\Omega_x$ be the exterior sites in its
radius-$r_*$ window, and set
\[
 (\Pi_u)_{yx}=\mathbf1_{y\in\Omega_x}
                        (A_{\Omega_x}+u)^{-1}_{yx},\qquad
 E_u=(A+u)\Pi_u-\mathrm{Id}.
\]
Exponential conjugation, followed by the resolvent identity, proves
in exponentially weighted row and column sums
\begin{equation}
 \|(A+u)^{-1}\|+\|\Pi_u\|\le\frac{C_L}{1+u},\qquad
 \|E_u\|\le\frac{C_Le^{-c_LM}}{1+u}.
 \label{rg:window-bounds}
\end{equation}
The residual lies outside its column window, so every contribution
crosses its radius.  The same proof for a difference of free histories
retains $\epsilon_h$.  These facts also hold if the exterior is empty.

Put
\[
 C_s=\tfrac12(\Pi_0+\Pi_0^*),\qquad
 \mu=-\Pi_0^*D^*F,\qquad r_s=D\mu+F.
\]
By \eqref{rg:window-bounds}, $C_s$ is positive for large $H$, and
it and all its principal marginals have upper and lower bounds at
fixed $L$.  The exact substitution $\xi=\mu+gZ$ uses the single
probability law $Z\sim N(0,C_s)$ and extracts
\begin{equation}
 (2\pi g^2)^{3n/2}\det(C_s)^{1/2}e^{-|r_s|^2/(2g^2)}
 \exp\left\{\tfrac12Z^*(C_s^{-1}-A)Z
                     -g^{-1}(A\mu+D^*F)^*Z\right\}.
 \label{rg:gaussian-ratio}
\end{equation}
Here $n$ is the number of exterior sites and the determinant includes
the three colors.  This identity follows by expanding the square;
the last exponential remains inside the expectation.

All nonlocal corrections in \eqref{rg:gaussian-ratio} are retained
as explicit paths.  With $E_s=AC_s-\mathrm{Id}$,
\begin{align}
 C_s^{-1}-A&=\sum_{r\ge1}(-E_s)^rA,\nonumber\\
 \log\det C_s&=-\operatorname{Tr}\log A+
           \sum_{r\ge1}\frac{(-1)^{r+1}}r\operatorname{Tr}E_s^r,
 \label{rg:ratio-series}\\
 (\log A)_{xx}&=\int_0^\infty
       [(1+u)^{-1}\mathrm{Id}-(A+u)^{-1}]_{xx}\,du,
 \quad (A+u)^{-1}=\Pi_u\sum_{r\ge0}(-E_u)^r.
 \nonumber
\end{align}
The leading-window term is a local scalar.  Extract its empty-pattern
value on every site, including frozen sites, and put the local
compensations into the adjacent core.  Residual chains have at least
one $E$ hop; their summed envelope, with any fixed requested support
exponent, is at most
\begin{equation}
 C_L\operatorname{poly}(M,r)
              e^{C B r}(C_Le^{-c_LM})^r.
 \label{rg:chain-price}
\end{equation}
The resolvent integral uses $(1+u)^{-2}$ for these terms.

On a stencil farther than $6M$ from the seeds, comparison with the
full empty-pattern operators gives
\begin{equation}
 r_s=O\bigl(C_LM^D(t^2+te^{-c_LM})\bigr),\quad
 |d(e^\mu T)|\le Ct/L+C_LM^Dt^2,\quad
 |\mu|\le Ct+C_LM^Dt^2.
 \label{rg:prediction}
\end{equation}
These estimates include normalized derivatives and free-history
differences.  To prove them, first replace the numerical kernels and
window inverse by the full empty-pattern kernels using
\eqref{rg:window-bounds}.  In the full calculation the first tangent
of $e^\mu T$ is exactly $P_hV$: the stack identities make the linear
residual vanish there, and $D^*D$ is invertible.  The remaining
terms are quadratic on the local $O(Mt)$ chart.  The first and
second fine-difference bounds for $P_h$ give the displayed constants
independent of $L$.  Everywhere, including near cores, the extended
formula still gives $|\mu|\le C_LM^Dt$.

\subsection{Prepared factors and a joint integral estimate}
\label{rg:prepared}

We describe the factors before expanding products.  This is necessary
because estimates on separate factors would not justify their products
under a common Gaussian.

The ownership rules and complete read sets are made explicit in
Sections~\ref{app:factor-ownership} and~\ref{app:read-sets}. A complete
read support includes geometric status queries as well as numerical
arguments; a ratio's Gaussian endpoint set is only part of that support.
Disjoint complete supports give the exact factorization of
Lemma~\ref{app:local-marginals}; overlapping factors use the joint
estimate below.

Mark, in each primary component, the selected input edges whose true
bad-bond inventory is supplied by an old covering label.  Require each
marked bit to be supplied once.  On every other selected input primary
edge retain $\mathbf1_{r_e\le t}$.  Nonprimary edges already have
$r_e\le t/2$ on the full product.  Thus these local rules reproduce
the old inventory constraint exactly.

For each primary component form a compulsory factor $B_c$ from its
assigned rows, with the selected affine squares subtracted.  Include
the negative stationary-square exponent within $8M$ of its seeds,
all primary tests and all defaults within $14M$, the frozen
normalization reciprocals, and the leading-window determinant
compensations.  Retain all true and contrary output-status tests
needed by these formulas.  In particular a component evaluated alone
does not infer the absence of a nearby output seed from unspecified
outside data.  Lemma~\ref{rg:core-reserve} gives
\begin{equation}
 |B_c|\le
 e^{-c_Lp^2s_c+C_LM^D(1+\log g^{-1})s_c}
                    \le e^{-c'_Lp^2s_c},
 \label{rg:core-bound}
\end{equation}
after the choices of $P_0,H$ made below.

On every exterior site put a smooth chart profile equal to one for
$|\xi|\le\pi/4$ and zero for $|\xi|\ge\pi/3$.  It preserves the
full sector integral and removes other angular preimages.  In addition,
attach a protector to every exterior spin read by a hard factor or by
$B_c$.  These are the protectors defined before Lemma~\ref{rg:core-reserve};
each remains attached to its factor and is never opened separately.  No extra hard predicate in the moving angle
$\xi=\log(qT^{-1})$ is imposed: alignment follows from the retained
sector profiles, or is imposed by a smooth larger plateau.  All merely
measurable dependencies are through physical input spins, fixed tags,
or output statuses that will not be field-differentiated.

The remaining factors are prepared as follows.  Each stated
exponential factor is then opened as $1+(e^V-1)$.
\begin{enumerate}
\item On unassigned $\mathcal F,\mathcal B$ rows use
$-b(|R|^2-|R^{\rm aff}|^2)/2$ with smooth angular cutoffs supported
in $|\xi|\le C_Lt$, whose plateaus contain all full-sector values.
Use smooth mean branches and no clipping.  The same preparation for
$-b\mathcal U$ is supported inside its smooth mean branch.
Their sup and fixed normalized derivative bounds are
$C_Lg\operatorname{poly}(M,p)$.

\item The unassigned stationary squares, necessarily farther than
$8M$ from the seeds, have bound $C_Lg^2\operatorname{poly}(M,p)$
by \eqref{rg:prediction}.  The unassigned $\mathcal C$ rows are zero.
The $\mathcal N,\mathcal Z$ rows have an exceptional bound
$C_Lg^{-D}e^{-c_LM}$, since $\mathcal N X'=0$ for full kernels.

\item A canonical term or regular error is called remote if its
entire projected graph is farther than $14M$ from the seeds.  For
a remote term replace its hard mask by a smooth product of graph-edge
profiles, one through $t/2$ and zero from $2t$.  The resulting
function is globally smooth: the error was defined through $4t$.
For a degree-$n$ canonical term with prefactor $b^r$ and coefficient
$a_X$, its bound is
\begin{equation}
 C_L|a_X|g^{n-2r}\operatorname{poly}(M,p,u).
 \label{rg:old-polynomial-bound}
\end{equation}
For an old error it is
$C_L\operatorname{poly}(M,p,s_X)[f_X]_{k,j}$.
For a nonremote term keep the actual mask and attach protectors.
The same undifferentiated bounds hold; it meets a compulsory factor
by the very test that made it nonremote.

\item Keep input and output kinetic truncation tails with both of
their paths and every mask query.  They have exceptional bounds
$C_Lg^{-D}e^{-c_LM}$ with their exponential path prices.  Input tails
may use the preceding smooth or protected hard prescription.  On an
output-empty extension use the analytic no-output-flag formula;
otherwise retain the actual output status.

\item Use the determinant chains and the linear and quadratic
Gaussian-ratio chains in \eqref{rg:ratio-series} literally.  A ratio
exponent $Q_\alpha$ has at most two Gaussian endpoints and
\begin{equation}
 |Q_\alpha(Z)|\le e_\alpha
                   \left(1+\sum_{x\in I_\alpha}|Z_x|^2\right).
 \label{rg:ratio-envelope}
\end{equation}
Even the eighth roots of $e_\alpha$ are summable with any of the
fixed support exponents required below: combine \eqref{rg:chain-price}
with the finite-range path count in Section~\ref{app:joint-application}.
The linear exponent uses
$-g^{-1}(A\mu+D^*F)$, which has the same exceptional accuracy.

\item Replace the exterior Haar Jacobian ratio by
$1+\chi(\xi)(J_{\exp}(\xi)/J_{\exp}(0)-1)$, where $\chi$ is a
smooth cutoff to $|\xi|\le C_Lt$ and has the required plateau.
Its second term has bound $C_Lg^2\operatorname{poly}(M,p)$ because
the Jacobian differs from its identity value quadratically.

\item Open the remaining edge and noise defaults, and the global
exterior chart profile, as one plus their signed failure letters.
For the remote noise default use
$\chi_0(|m|)G_{\mathrm{noise}}(V-m/|m|)$, where $\chi_0=1$ above
$3/4$ and $\chi_0=0$ below $2/3$; evaluate the normalized mean only
on this smooth branch.  This equals the original default on the
full sector.  The failure and all its nonzero derivative supports
force a large Gaussian coordinate, as proved below.

\item An old covering weight retains its own sup envelope and all
its protectors.  Every one of its inventory edges is a selected
input primary, so the label meets a compulsory factor.  Its
measurable function is never differentiated.
\end{enumerate}
The empty-pattern leading determinant, the Gaussian scalar powers,
the Haar constant and the observation denominator have already been
extracted.  Their local changes were put into $B_c$.  In particular
there is no unaccounted ordinary factor of order zero.

Every preparation above preserves the original full sector product:
the required cutoffs are one and the smooth branches agree there.
Only after these preparations do we open factors.  A resulting term
need not satisfy every original default, and its primary geometry is
not recomputed from the values of $Z$.

\begin{lemma}[Joint prepared-factor bounds]
\label{rg:joint}
The prepared factors have nonnegative majorants $a_\alpha$ such that,
for every admissible finite selection $A$ in a fixed primary pattern,
including its compulsory component factors, and every
allowed normalized derivative of total order $r\le k$,
\begin{equation}
 \left|D^r\mathbb E\prod_{\alpha\in A}F_\alpha\right|
 \le C_k\left(1+\sum_{\alpha\in A}s_\alpha\right)^{d_k}
                      \prod_{\alpha\in A}a_\alpha.
 \label{rg:joint-product}
\end{equation}
For nonempty output inventory only the undifferentiated and parameter
comparison forms are asserted. A single marked discrepancy retains its
own size in the corresponding product bound. The load polynomial is
absorbed by a reserved amount of support exponent. The product estimate
is proved explicitly in Appendix~\ref{supp:joint-majorants}, with the
concrete eight-class verification at its end. The majorants have the
following summed bounds.
An ordinary primitive of order $r>0$ has a summed majorant
$C_Lg^r\operatorname{poly}(M,p)$, with any fixed required support
exponent.  The old errors retain their actual cap $\delta_j$.
Exceptional primitives retain an arbitrary fixed power of $g$,
and an output-seeded connected component retains a bound stronger
than $e^{-p^{3/10}}$.

For a component with empty output inventory, the same statements
hold through $k$ normalized field derivatives on its
no-output-flag extension, and through $k-1$ derivatives for a single
input or parameter discrepancy.  For a component with nonempty
output inventory, only sup bounds and coupling/history differences
at fixed $V$ and fixed output statuses are asserted.  All
discrepancies retain their own input difference or the factor
$\epsilon_h+\lambda_g$, where $\lambda_g=|b_2-b_1|/H_j$.
\end{lemma}
\begin{proof}
For the first prepared rows, $R-R^{\rm aff}=O(C_LM^Dt^2)$ and
$R+R^{\rm aff}=O(C_LM^Dt)$ on their cutoffs.  Their product times
$b$ is $C_Lg\operatorname{poly}(M,p)$.  The mean error in
$\mathcal U$ is quadratic and is multiplied by an output gradient.
Each normalized derivative replaces an angular factor by a direction
of the same size.  A cutoff derivative costs $t^{-1}$ and is paired
with such a direction.  Thus the orders are preserved.  Substitution
of $\mu+gZ$ adds fixed Gaussian polynomials and fixed $M,p$ losses;
it does not change the order.  The same argument, using the
quadratic jet of $r_s$, proves the stationary and Jacobian bounds.

For \eqref{rg:old-polynomial-bound}, each relative slot is bounded
by $2t$ times the length of its path on the smooth graph cutoff.
The old coefficient exponential pays every fixed power of that
length, including the compensation coefficients.  An old error is
multiplied by a smooth cutoff supported strictly inside its $C^k$
domain and extended by zero.  Its derivatives use its prescribed
norm, with only fixed powers of graph size and translated direction
weights.  This argument does not require a large graph to fit in one
angular chart.  For a map discrepancy acting on an unchanged error,
one extra derivative uses order at most $k$; a difference of errors
itself is used only through order $k-1$.

For ratio letters, an elementary exponential estimate gives
\[
 |e^{Q_\alpha}-1|
 \le C\sqrt{e_\alpha}
     \exp\left\{C\sqrt{e_\alpha}
                (1+\sum_{I_\alpha}|Z_x|^2)\right\}.
\]
All endpoint loads are small in per-site sum.  For any fixed moment
exponent, the Gaussian determinant formula therefore bounds a
product over distinct ratio letters by
$C_L^{n}\prod_\alpha e_\alpha^{1/4}$ in that moment norm.
Indeed the matrix in the quadratic exponential has sufficiently
small norm, and its logarithmic determinant is bounded by a constant
times its trace.  The eighth-root reserve pays fixed differentiated
Gaussian polynomials as well.  No independence of overlapping
letters is used.

For a remote edge failure, \eqref{rg:prediction} and
$\max|Z_x|<c_Lp$ on its stencil would put the chord strictly below
$t/4$.  For a remote noise failure they would give $|m|>3/4$ and
$|V-m/|m||<Wt$, after $W$ is chosen large.  The same conclusions
hold on every derivative support.  A chart failure has
$|\xi|\ge\pi/4$; the global bound $|\mu|\le C_LM^Dt=o(1)$ gives
the stronger condition $|Z_x|\ge c/g$.
Among $n_d$ distinct local failure tests, select at least
$c_Ln_d/M^D$ with disjoint stencils.  This follows from the bounded
polynomial number of test types and placements meeting a stencil.
Choose an offending site in each.  Exponential Chebyshev with a
fixed small positive quadratic coefficient and the bounded
covariance gives
\begin{equation}
 \mathbb P(\text{all specified failures})
             \le e^{-c_Lp^2n_d/M^D}.
 \label{rg:joint-rarity}
\end{equation}
The number of offending-site choices is absorbed into this bound.
Again the Gaussian sites need not be independent.  A fixed Holder
inequality combines \eqref{rg:joint-rarity} with the ratio estimate.

It remains to justify derivatives through hard inputs.  Take the
one marginal on the union $I$ of all read Gaussian coordinates.
On each protected exterior spin solve
\[
 \partial_vq_x+W_x\cdot\nabla_{Z_x}q_x=0.
\]
The differential of the exponential chart is invertible on the
protector's compact support.  Its inverse costs $g^{-1}$; extend
$W_x$ smoothly outside a slightly larger chart and set other
components to zero.  All hard physical-spin arguments are then
fixed simultaneously.  Integration by parts for this full marginal
gives
\begin{equation}
 \partial_v\mathbb E F=\mathbb E\left[
  (\partial_v+W\cdot\nabla_Z)F+
  F\{\operatorname{div}W-Z[I]^*C_s[I,I]^{-1}W[I]\}\right],
 \label{rg:transport}
\end{equation}
with the additional covariance score when the covariance varies.
Frozen spins and tags are fixed in this operation.  Explicit angle
dependence in the affine squares, smooth frozen-angle extensions,
chart profiles and protectors is differentiated normally.  There
is no hard moving-angle predicate to which
\eqref{rg:transport} would be inapplicable.

Repeating a fixed number of times costs only a fixed polynomial in
$g^{-1},M,p,|I|$ and the Gaussian variables.  Every connected hard
use has a retained reserve: an old covering label supplies a
primary bit; a nonremote mask meets a primary by its definition;
a voluntarily hard tail has its own exceptional accuracy.
The core bound also holds on differentiated profile supports by
Lemma~\ref{rg:core-reserve}.  For a parameter difference interpolate
the two core exponents at the same physical $q$; both endpoints
satisfy that reserve, hence so does their interpolation.

Only finitely many factors are differentiated at fixed total order.
Their placement count and their Gaussian moment cost are polynomial
in the total load, not exponential in the number of factors.
All undifferentiated ordinary factors have deterministic sup bounds.
Thus a single fixed Holder exponent suffices, independent of that
number.  The exceptional reserves absorb the finite negative powers
of $g$.  The restriction on output field derivatives in the statement
is essential: the actual output-status indicators are not smooth.

Finally convert the individual bounds to support-hit sums.  An old
anchor contributing to a new hit lies within $CM'(1+s/L)$ of it;
there are only polynomially many such positions.  For a fixed
requested exponent $B$, \eqref{rg:projection} gives
\[
 e^{Bs'-As}(1+s)^D
 \le e^{BD_*}(1+s)^De^{-(A-BD_*/L)s}.
\]
Choose $L$ so that $BD_*/L<A/4$.  The unused exponential pays these
polynomials.  Canonical paths have their stronger analytic length
weight; take $H$ so that a mesoscopic price consumes less than a
fixed fraction of it.  Core patterns have at most
$(C_LM^D)^{Cs_c}$ choices at a support hit, paid by
\eqref{rg:core-bound}.  Chain sums follow from
\eqref{rg:chain-price}.  This proves all stated summed bounds.
\end{proof}

\subsection{Connected sums and exact reassembly}
\label{rg:connected}

Here is the elementary estimate used for each connected sum.  Suppose
primitive majorants $a_\alpha$ on supports $S_\alpha$ satisfy
\begin{equation}
 \sup_x\sum_{\alpha:x\in S_\alpha}
            a_\alpha e^{(B+2)s_\alpha}\le\eta,
 \qquad |S_\alpha|\le C_LM^2s_\alpha,
 \qquad C_LM^2\eta\le\tfrac14.
 \label{rg:tree-condition}
\end{equation}
The sum of the rooted trees above a root of load $s$ is at most
$e^s$: assuming this bound for the descendants, the unordered
children have exponential sum at most $\exp(C_LM^2\eta s)$.
Induction on tree depth and monotone convergence prove the assertion.
Connected sets are bounded by their spanning trees.  For a hard-core
Mayer logarithm the absolute connected-graph coefficient is at most
the number of spanning trees.  To see this, fix a total edge order
and partition connected graphs according to their minimal spanning
tree.  Each class is an interval of optional extra edges; its
alternating edge sum is zero or has absolute value one.  Treat
repeated labels as distinct numbered vertices before dividing by
the factorial.  This gives the same rooted-tree estimate.

The complete supports make the algebra match these bounds.  Give
each noncore read position its $20M$ halo.  A mean or covariance
window reads only this halo, including the $3M$ hole-status test.
The mask radius $O(\log(1/t))$ is smaller than $r_*$ for large $H$.
All reads of a core, including its $14M$ defaults and their windows,
fit inside the $50M$ footprint with these margins.  Long determinant
or inverse chains retain each successive window as a separate
charged hop.  Two disjoint complete supports therefore have zero
cross covariance under $C_s$ and disjoint frozen/tag variables.
Their individual marginal entries are unchanged when outside
primaries are deleted.  Their integrals consequently factor exactly.

Group each connected collection of selected prepared factors into
an integrated polymer, retaining its original compatibility,
coverage and status checks.  A disjoint family of such polymers
with the full required output inventory reconstructs exactly the
original selection of prepared factors.  An output-empty polymer
can be required to have all of its output graph good: otherwise the
mandatory cover of a bad edge in that graph would intersect it and
belong to the same component.  It is therefore a regular masked
function with the extension furnished by Lemma~\ref{rg:joint}.

Apply the Mayer expansion only to these output-empty polymers and
write its logarithm as $J=\sum_i1_{X_i}j_i$.  For a fixed compatible
family of output-seeded polymers, the allowed background has
logarithm equal to $J$ minus just those Mayer terms whose supports
meet that family.  Factor out $e^J$, open each of these forbidden
terms once as an exponential-minus-one letter with the opposite
sign, and join it to the seeded polymers it meets.  Every new
component has nonempty output inventory.  We obtain exactly
\begin{equation}
 e^{v_{\rm step}|\Lambda'|-b\mathcal E'}e^J\Xi'_{\rm pre},
 \label{rg:preliminary-output}
\end{equation}
where $\Xi'_{\rm pre}$ is again a mandatory covering gas.

The same operation implements every later mask strengthening.  The
logarithmic correction $(1_X-1_Y)j$ for a stronger mask $Y$ can be
nonzero only if its support contains a true output bad edge.  In a
full family it therefore meets that edge's covering label.  Open
the correction once and regroup.  If supports have been enlarged,
regroup by the enlarged supports while retaining the original
compatibilities inside the new weights.  Connected-component
decomposition is a bijection, so these operations preserve the
density and the inventory exactly.

For countable lists, do the algebra first for finite subsets of the
prepared factors.  The resulting identity is for that truncated
prepared expression.  Lemma~\ref{rg:joint} and
\eqref{rg:tree-condition} give a volume-exponential absolute
majorant, also for the required derivatives.  Dominated convergence
then gives \eqref{rg:preliminary-output} for the original full
expression.  No positivity of the truncated signed gases is used.

\subsection{The canonical coefficients}
\label{rg:canonical}

The canonical map is a finite-order algebraic operation on the full
bulk kernels.  It is useful to define it before comparing it with
the exact integration.  Use the empty primary pattern, no masks or
cutoffs, and the full free kernels.  Write
\[
 R=(X-\nu,Y-w),\qquad(\nu,w)=\mathcal M X',\qquad
 U_E=X'_E\cdot[\mathcal M^*(0,\sqrt a(m-m/|m|))]_E.
\]
The full energy identity is
$\mathcal E(q)+|Y|^2/2-\mathcal E'(V)=|R|^2/2+\sum_EU_E$.
Use its affine row $R^{\rm aff}=D_f\xi+F_f$, its mean
$\mu_f=-C_hD_f^*F_f$, and its stationary residual
$R_{\mathrm{stat}}=D_f\mu_f+F_f$.  Substitute $\xi=\mu_f+gZ$,
$Z\sim N(0,C_h)$, in the following negative-log vertices:
\begin{equation}
 \frac b2(|R_i|^2-|R_i^{\rm aff}|^2),\qquad
 bU_E,\qquad\frac b2|R_{\mathrm{stat},i}|^2,\qquad
 \text{the old canonical slots on }q,\qquad
 -\log\frac{J_{\exp}(\xi_x)}{J_{\exp}(0)}.
 \label{rg:canonical-vertices}
\end{equation}
An anchor is removed by common right multiplication before using
its coordinates $y_o$.  Count $y_o$ and $gZ$ as order $g$, and
retain total order at most four in the negative logarithm of the
Gaussian expectation.  Equivalently, for each ordered multiset of
$v\le4$ vertices use its connected Wick expectation with coefficient
$(-1)^{v+1}/v!$.  Both internal and intervertex Wick pairs are included.
The coefficient sums converge absolutely by the free exponential
bounds and \eqref{rg:canonical-norm}.

Write the resulting expression as
$\sum_{r,d}g^{-2d}F_{r,d}(y)$, with $r\le4$.
Thus $F_{r,d}$ is the coefficient polynomial without its displayed
power of $g$, and is homogeneous of field degree $r+2d$.
Scalar polynomials are extracted per unit volume.  Let
$S_{K',o}^{[n]}$ be the degree-$n$ part, in the relative $y_o$ chart,
of the analytic unmasked output kinetic density, assigned to its
anchor by splitting edge and row contributions between their
endpoints.  Its affine quadratic Hessian is one.  Its odd-degree
parts vanish, since it is quadratic in ambient spins and every dot
product has even total degree in this chart.

Let $\operatorname{Comp}$ denote the orbitwise quadratic
compensation described after \eqref{rg:canonical-norm}.  The map is
\begin{align}
 P'_4&=F_{2,1},&
 \alpha&=\text{affine Hessian of }F_{2,0},&
 I'_2&=\operatorname{Comp}(F_{2,0}-\alpha S_{K'}^{[2]}),
 \nonumber\\
 P'_5&=F_{3,1},& I'_3&=F_{3,0},\nonumber\\
 P'_6&=F_{4,1},&
 I'_4&=F_{4,0}-\alpha P'_4-\alpha S_{K'}^{[4]},&
 \kappa&=\text{affine Hessian of }F_{4,-1},
 \nonumber\\
 J'_2&=\operatorname{Comp}(F_{4,-1}-\kappa S_{K'}^{[2]}).
 \label{rg:canonical-map}
\end{align}
Here and below the affine Hessian means its bulk sum per anchor.
After the bulk affine coefficient has been removed, the sum of the
individual orbit compensation coefficients is zero; thus
$\operatorname{Comp}$ does not change that polynomial sum.

Appendix~\ref{supp:canonical-fixed-norm} specifies one common coefficient
norm and proves the diagonal estimate, including compensation and the
translation-invariant bulk symmetry needed by the quadratic slots.

\begin{lemma}[Completeness and contraction of the canonical map]
\label{rg:canonical-contraction}
The list \eqref{rg:slots} contains every nonscalar canonical term
through total order four.  In that order the map
\eqref{rg:canonical-map} is triangular, with same-slot linear
response at most $C/L$ in the coefficient norm, where $C$ is fixed
before $L$.  All other responses are bounded by constants depending
on $L$ and the preceding caps.  The kick $\alpha$ does not depend
on old $I_2$, and $\kappa$ does not depend on old $J_2$ at their
respective orders.

Consequently the caps in \eqref{rg:canonical-norm} can be chosen
successively so that the canonical orbit starting from zero remains
bounded.  This orbit, and its kicks $\alpha_h,\kappa_h$, approach
limits exponentially in $h$.
\end{lemma}
\begin{proof}
The first vertex of \eqref{rg:canonical-vertices} begins with a cubic,
of total order one.  Its cubic contains a Gaussian field: at $Z=0$
both $R$ and $R^{\rm aff}$ have zero first tangent by the harmonic
identity.  The $U$ vertex begins in degree four.  Its possible cubic
would pair the real degree-two part of $m-m/|m|$ with the imaginary
first-degree part of $X'$, and is zero.  The stationary-square
vertex begins in degree four, and the Haar vertex in degree two.

A connected Wick graph on $v$ vertices requires at least $v-1$
intervertex pairs.  Each vertex carries at most $g^{-2}$ and each
pair supplies $g^2$, so its remaining prefactor is at most $g^{-2}$.
Write it as $g^{-2d}$, with the integer $d\le1$.  A linear field is
forbidden because the stabilizer $SO(3)$ has no invariant vector.
The possible degree-three term at order one is zero by the
Gaussian-field observation above.  At orders two, three and four,
the remaining possibilities are respectively
\[
 (bP_4,I_2,\text{scalar}),\qquad
 (bP_5,I_3),\qquad
 (bP_6,I_4,b^{-1}J_2,\text{scalar}).
\]
The full invariant-tensor basis includes every color contraction,
including odd alternating tensors.  No additional tensor has been
discarded by a parity assumption.

For a quadratic form, conjugation invariance gives a color dot
product.  On affine fields it has the form
$\sum_{\mu,\nu}T_{\mu\nu}a_\mu\cdot a_\nu$.
Spatial reflections give $T_{12}=0$ and the quarter turn gives
$T_{11}=T_{22}$.  Thus there is exactly one affine quadratic
coefficient.  Its removal is precisely the stiffness extraction
in \eqref{rg:canonical-map}; polarization then gives vanishing on
every pair of affine fields.  A one-site invariant potential is
constant, so there is no additional relevant potential term.

Composition cannot lower the total order of an old slot, since an
anchor-relative field has no constant term.  At its own order an
old slot therefore enters only through a singleton Gaussian
expectation and its linear field map.  Contractions lower $d$ and
cannot raise it.  Ordering equal total orders by decreasing $d$
gives exactly the triangular order \eqref{rg:slots}.  The
subtractions in \eqref{rg:canonical-map} preserve it.  In particular
the term $\alpha P'_4$ is necessary to express the output with
coupling $b+\alpha+\kappa/b$.

For its diagonal response, a relative field at $x$ is replaced by
$\sum_z(P_h(x,z)-P_h(o,z))y_o(z)$.  The free fine-difference bounds
give, with any of the fixed smaller exponential weights needed here,
\begin{equation}
 \sum_z e^{c\sigma|z-[o]|}|P_h(x,z)-P_h(o,z)|
 \le \frac{C|x-o|}{L}
                     e^{C_1\sigma(1+|x-o|/L)}.
 \label{rg:row-difference}
\end{equation}
Its constants are independent of $L$.  A degree-$n$ substitution
therefore has a factor $L^{-n}$; the $L^2$ old anchors over a new
anchor give $L^{2-n}\le L^{-1}$ for $n\ge3$.  Fixed powers of old
path lengths are absorbed by the unused old exponential, uniformly
once $L\ge2C_1$.

For a compensated quadratic write this row difference as
\[
 (x-o)\cdot\nabla_fP_h(o,\cdot)+R_x,
 \qquad
 \|R_x\|_{\exp}\le
        \frac{C|x-o|^2}{L^2}e^{C_1\sigma(1+|x-o|/L)}.
\]
The product of the two affine pieces cancels on the whole old
compensated orbit.  The remaining coefficient polynomials have a
factor $L^{-3}$, which gives $C/L$ after the anchor count.  This
cancellation is performed before taking absolute coefficient norms.
Output compensation adds only fixed polynomial path factors; spare
exponential width pays them with a constant fixed before $L$.
Moreover the aggregate transferred affine Hessian is exactly zero,
because $P_h$ preserves affine fields.  This proves the assertions
about $I_2,\alpha$ and $J_2,\kappa$.

All nondiagonal terms contain a bounded number of vertices.  Their
connected position sums are bounded by exponentially summable free
kernels and old coefficient lists, so they have finite $C_L$ bounds,
also for differences.  Choose $L$ for diagonal contraction and then
the caps successively.  Positive triangular weights give a common
contraction factor $q_c<1$ for differences in a bounded orbit.
The free-history forcing is $C_L L^{-h}$.  Comparing an orbit of
length $h$ with any longer one aligned for its last $h$ steps gives
\[
 d_h\le q_c^h C_L+
             C_L\sum_{r=0}^{h-1}q_c^{h-1-r}L^{-r}
                \le C_L\rho_c^h
\]
for a fixed $\max(q_c,L^{-1})<\rho_c<1$.
This is uniform in the extra history, hence proves the Cauchy and
kick assertions.  No interacting fixed point has been assumed.
\end{proof}

\subsection{The regular error and the exact Taylor comparison}
\label{rg:error}

The one old-error contribution which cannot be discarded by order
counting is its isolated linear response
\begin{equation}
 \mathcal T f_X(V)=\mathbb E_{C_s}
       [\chi_X(q)f_X(q)],\qquad
 q_x=\exp(\mu_x+gZ_x)V_{[x]},
 \label{rg:error-transfer}
\end{equation}
in the empty primary pattern.  Give it the projected old anchor.
For $s_X\le L^{1/4}$ include a complete output good-field ball of
radius $CM'$ and its required stencils.  Increasing $C_{\mathrm{adm}}$
makes these cases nonwinding.  Any strengthening of an existing
mask is compensated by the exact operation following
\eqref{rg:preliminary-output}.

\begin{lemma}[Contraction of old regular errors]
\label{rg:error-contraction}
For each fixed required output exponent $B$, and sufficiently large
$L$ followed by $H$,
\begin{equation}
 \|\mathcal T f\|_{k,j-1,B}
                  \le (C/L+o(1))\|f\|_{k,j,A}.
 \label{rg:error-contraction-bound}
\end{equation}
The same bound applies to an input difference through order $k-1$.
A change of map acting on an unchanged input instead costs
$C_L\operatorname{poly}(M,p)\delta_j(\epsilon_h+\lambda_g)$
through order $k-1$.
\end{lemma}
\begin{proof}
Retain an even smooth Gaussian guard $|Z_x|\le c_Lp$ on all read
sites.  Its complement costs
$C_LM^D(1+s_X)^De^{-c_Lp^2}[f_X]_{k,j}$, including the derivatives
at issue, by the union tail estimate and Lemma~\ref{rg:joint}.
On the guard the old cutoff $\chi_X$ is one, also when the
fluctuation is multiplied by any number in $[-1,1]$.

For output orders at most two, Taylor expansion twice in that
fluctuation removes it with error
\begin{equation}
 C_LM^D(1+s_X)^Dp^{-2}[f_X]_{k,j}.
 \label{rg:fluctuation-omission}
\end{equation}
The first term integrates to zero by the even guard.  Each
fluctuation insertion has size at most $C\sup|Z|/p$ in the old
direction norm; mixed map derivatives retain these two factors.
Only four old derivatives are needed.  Choose $P_0$ larger than
the fixed $M$-power requirement, so this is a delayed small factor.

For a short label remove the anchor spin by right multiplication
and write
\[
 A_x=\log(q_x^0(q_o^0)^{-1}),\quad q_x^0=e^{\mu_x}V_{[x]},
 \qquad
 \mathfrak b_x=(x-o)\cdot\nabla_fP_h(o,\cdot)y.
\]
In the old normalized direction norm, including the low output
derivatives,
\begin{equation}
 \|A\|\le C/L+o(1),\qquad
 \|A-\mathfrak b\|\le C/L^2+o(1).
 \label{rg:affine-prediction}
\end{equation}
For $r=|x-o|\le L$, the first and second fine-difference bounds
give respectively $Ct'r/L$ and $Ct'r^2/L^2$.  Dividing by the old
direction bound $t(1+r)^8$ gives the two displayed gains, since
$t'/t$ is bounded.  For $r>L$, use instead the bounded row of $P_h$:
the ratio of direction weights is at most
$C(1+r/L)^8/(1+r)^8\le CL^{-8}$.
The affine term divided by the old weight is at most
$Cr/[L(1+r)^8]\le CL^{-8}$ as well.  This treats distant sites
without extending a small-displacement Taylor bound beyond its range.
The same estimates apply to the linear parts of the first two output
derivatives.  Nonlinear chart terms have an extra $t$ with fixed
polynomial $M$ costs, which tend to zero after $L$ is fixed.
The Gaussian guard
permits the same first-derivative bound for the relative-angle
map with fluctuation retained, while every higher map derivative
is $o(1)$ in this norm.

To see the cancellation also in the differentiated norms, put
$F_X(A)=f_X(e^A)$ and $\mathcal B_X=D^2F_X(0)$ after removing the
anchor spin.  The normalization of $f_X$ gives
$F_X(0)=DF_X(0)=0$, and $\mathcal B_X$ vanishes on every pair of
affine tangent fields.  For two normalized output directions write
$A_i=\partial_iA$ and $A_{12}=\partial_1\partial_2A$.  Taylor's
formula and the chain rule give
\begin{align*}
 F_X(A)&=\tfrac12\mathcal B_X(A,A)+R_0,\\
 \partial_iF_X(A)&=\mathcal B_X(A,A_i)+R_i,\\
 \partial_1\partial_2F_X(A)
   &=\mathcal B_X(A_1,A_2)+\mathcal B_X(A,A_{12})+R_{12}.
\end{align*}
The bounds \eqref{rg:affine-prediction}, including these derivatives,
give
\[
 |R_0|+|R_1|+|R_2|+|R_{12}|
       \le \bigl(C/L^3+o(1)\bigr)[f_X]_{k,j}.
\]
For example the second-derivative remainder is bounded by a constant
times
$(\|A\|\|A_1\|\|A_2\|+\|A\|^2\|A_{12}\|)[f_X]_{k,j}$;
only three old derivatives are needed here.  Every derivative of
$\mathfrak b_x=(x-o)\cdot\nabla_fP_h(o,\cdot)y$ remains affine in
$x-o$.  Replacing each $A$ or its derivatives in the bilinear terms
by the corresponding $\mathfrak b$ therefore gives zero.  Each
remaining term contains one difference controlled by $C/L^2+o(1)$
and one factor controlled by $C/L+o(1)$, proving the same gain for
the bilinear terms.  The chart conversion costs already included above
and the fluctuation-omission error tend to zero with $H$.

The Taylor segment stays in the old graph domain because
\eqref{rg:prediction} puts its chords strictly inside the cutoff
plateau.  For output orders three through $k$, retain the fluctuation
and apply the chain rule directly on the Gaussian guard.  A term has
either at least three first map derivatives, giving $C/L^3+o(1)$,
or a higher map derivative, giving $o(1)$ for fixed $L$.
This uses no old derivative above the requested order; it does not
differentiate a Taylor remainder through order $k$.

There are $L^2$ old anchors over a new anchor.  On these short
labels the output load is bounded independently of $L$, so its
fixed exponential price gives $C/L+o(1)$.  On labels
$s_X>L^{1/4}$, including winding ones, use the direct composition
bound $C(1+s_X)^D[f_X]_{k,j}$.  There is no loss of a bare $M$
power in its leading first-direction comparison: the absolute
first jet is the bounded row of $P_h$.  Higher map derivatives
have extra $t$ factors paying their $M$ powers.  The unused old
size exponential in \eqref{rg:projection} gives
$e^{-cL^{1/4}}$, which pays $L^2$ and the displayed polynomial.
This proves \eqref{rg:error-contraction-bound}.

An input difference uses exactly the same argument in order
$k-1$.  For a changed map, ordinary composition estimates use one
extra derivative of the unchanged error, at most $k$, and give the
stated map-difference bound.  Thus iteration does not lose one
derivative at each scale: every individual output is still $C^k$,
and every discrepancy is measured in $C^{k-1}$.
\end{proof}

We now compare the exact logarithm in
\eqref{rg:preliminary-output} with \eqref{rg:canonical-map}.
An ordinary positive-order primitive of order $r$ has majorant
$C_Lg^r\operatorname{poly}(M,p)$.  Inflate it by $g^{-.96r}$ in
\eqref{rg:tree-condition}.  Its remaining $g^{.04r}$ still beats
every fixed logarithmic loss and the factor $M^2$.  The connected
sum of terms of total order at least five therefore retains
$g^{4.8}$.  A smaller exponent pays the remaining fixed derivative
and logarithmic costs.  Thus, after removing the isolated old-error
response, all nonprincipal terms cost
\begin{equation}
 C_Lg^{4.6},\qquad C_Lg^{4.5}\tau\text{ for discrepancies},
 \quad \tau=u+\epsilon_h+\lambda_g,
 \label{rg:high-order-remainder}
\end{equation}
in the needed anchored norms and a sufficiently large fixed support
exponent.  Old errors are counted as order four here but retain
their stronger actual cap; every nonisolated use has total order
at least five.  Exceptional terms and canonical paths longer than
$M$ have arbitrary power accuracy.

For the principal terms, group the at most four origins at one
anchor and strengthen their masks to a common ball.  Choose
$C_\chi$ large and then $L$ large: the old canonical balls shrink
by $L$, while the fixed number of new stencils uses only fixed
multiples of $M$.  The common ball contains every required read.
Its complement is compensated in the covering gas.  On that ball
and the even Gaussian guard, every prepared factor has its smooth
vacuum formula.  Taylor expansion through total order four, then
restoration of full Gaussian polynomial moments, costs
$C_Lg^5\operatorname{poly}(M,p)$, also in the normalized derivative
norms.  The largest vertex degree required is six; a higher
normalized derivative retains the same power because each
direction supplies a factor $t$.

On each fixed origin multiset the exact product and Mayer-log
identities are identities of these finite polynomials.  They give
precisely its connected Wick expectation under the short Gaussian,
including repeated origins and the Jacobian logarithm.  Replacing
short kernels and covariance by full ones costs exponential
accuracy in $M$, with all fixed position moments and discrepancies.
Any connection that appears only after this replacement contains
a long kernel path and has that same accuracy.  The principal
answer is therefore exactly \eqref{rg:canonical-map}, up to
\eqref{rg:high-order-remainder}.  This proves that the formal
coefficient operation describes the actual integration to the
claimed accuracy.

\subsection{Normalization, periods, and closure}
\label{rg:normalization}

The remaining task is to remove the constant and affine-Hessian parts
of the regular error without changing the density. These two extractions
determine the scalar increment and the remaining kinetic increment
$\Delta$.

Put $b^-=b+\alpha+\kappa/b$.  Inserting the full canonical lists
with their masks and expressing their displayed prefactors in terms
of $b^-$ leaves a raw regular list.  Denote its norm by $R$ and the
norm of its difference for two inputs by $R^\Delta$.  The preceding
two subsections give
\begin{align}
 R&\le q_1\delta_j+C_Lg^{4.6},\nonumber\\
 R^\Delta&\le q_1\delta_j^\Delta+
 C_L\operatorname{poly}(M,p)\delta_j(\epsilon_h+\lambda_g)
       +C_Lg^{4.5}\tau,
 \qquad q_1=C/L+o(1).
 \label{rg:raw-error}
\end{align}
Here $\delta_j^\Delta$ is the actual input error discrepancy, before
division by $\delta_j$, and $\lambda_g=|b_2-b_1|/H_j$ up to a fixed
constant.  The constant multiplying $L^{-1}$ is fixed before $L$;
the $o(1)$ is made small by the subsequent choice of $H$.
These estimates hold with any of the finitely many larger support
exponents reserved for the operations that follow.

We describe the insertion of the kinetic terms explicitly.  Assign
each analytic base path to an anchor, using the convention in
\eqref{rg:canonical-map}, and require on its mask every edge queried
by the relevant $m'_e$.  On that mask it is the analytic chord
expression.  Its Taylor remainder after degree four for the
$\alpha$ term, or degree two for the $\kappa/b$ term, has the
accuracy in \eqref{rg:high-order-remainder}.  A long path pays its
position and derivative powers from exponential localization.
Off the mask keep the actual clipped path expression as a covering
correction.  Its site-hit bound is a small power of $g$: the direct
clipped energy is at most $Ct'$, and an active second-slot row is
bounded using $m'_e=1$ and its exponentially weighted chord bound.
Thus the exact insertion uses no Taylor formula outside its domain.

For a nonwinding raw error at anchor $o$, let $v_i=f_i(1)$ and let
$b_i$ be its unit-affine Hessian.  The direction norm gives
\[
 |v_i|\le[f_i]_{k,j-1},\qquad
 |b_i|\le C(t')^{-2}[f_i]_{k,j-1}.
\]
There is no support-size loss in the second estimate: a unit
affine field at distance $r$ is bounded by a constant times
$(1+r)^8$.  After strengthening the mask to include the nearest
neighbors of $o$, use the exact identity
\begin{equation}
 1_Xf_i=v_i+b_iS_o+
        1_X(f_i-v_i-b_iS_o)
       -(1-1_X)(v_i+b_iS_o).
 \label{rg:normalization-identity}
\end{equation}
The middle term has zero value and zero affine Hessian.  On its
extension, $S_o$ and its normalized derivatives are bounded by
$C(t')^2$, so its norm is at most a fixed multiple of that of
$f_i$.  The last term meets an output bad bond and is included in
the covering gas by the exact regrouping in
Section~\ref{rg:connected}.

Sum the bulk $v_i$ per anchor into the volume scalar and put
$\Delta=\sum_i b_i$, also using the bulk lists.  Replacing $b^-$
by $b'=b^-+\Delta$ leaves the negative-log correction
\begin{equation}
              -\Delta\bigl(\mathcal E'-\sum_oS_o\bigr).
 \label{rg:kinetic-reset}
\end{equation}
Allocate the compensating $S_o$ to each base path in proportion
to that path's affine Hessian.  These Hessians sum to one.  On its
path mask each resulting expression has zero affine Hessian,
and the sum of its regular norms is at most
$C(t')^2|\Delta|$.  The constant here is independent of $L$:
it uses only the uniform exponentially weighted free-row moments
and the fixed direction order.  Mask complements are again
covering corrections.  The changes in the seven displayed
canonical prefactors are smaller; for example the change
$\Delta P'_4$ has an additional factor $(t')^4$.
Consequently the final regular norm and coupling error satisfy
\begin{equation}
 \|f'\|_{k,j-1,A}\le Cq_1\delta_j+C_Lg^{4.6},
 \qquad
 |\Delta|\le C_LH_j^{-1.05}.
 \label{rg:closed-error}
\end{equation}
The same operations give the difference estimate in
\eqref{rg:raw-error}, up to fixed factors and the indicated
$(t')^{-2}$ for $\Delta$.  In particular no derivative above
order $k$ of an input error has been used.

We next justify using bulk $v_i,b_i$ on every admitted period.
Every finite-period free entry is the fold of its lifted paths.
All primitive choices retain complete lifted read records.
A connected calculation with load below $c_*n'/M'$ therefore
agrees exactly with its bulk calculation.  A disagreement has
load at least $c_*n'/M'$, and hence has an additional factor
\[
                 \exp(-c_*A_1n'/M')
\]
from a reserved support exponent $A_1$.  Keep a period-dependent
constant or Hessian disagreement as an actual winding error,
with a declared record of load $O(n'/M')$.  The extra exponential
pays this record.  Its constant is an allowed winding term; its
Hessian is multiplied by $S_o$, recovering the factor $(t')^2$
in the regular norm.  An off-mask correction has the same
covering interpretation as before.  Already winding inputs
remain winding because $s'\ge4s/L$ in \eqref{rg:projection}.
Thus all period discrepancies stay in the error or covering
lists.  The extracted coupling and volume scalar are precisely
the bulk ones and are independent of the period.

Here is a finite order of choices closing all the estimates.
Reserve a sufficiently large fixed support exponent $E_*$ for
the finitely many tree, mask, orbit and winding operations, and
let $D_*$ bound their geometric dilations.  First choose $L$
large enough that
\[
 E_*D_*/L<A/4,\qquad
 C_{\rm lead}e^{2E_*D_*}/L<1/16,
\]
as well as the free and core requirements already stated.
$C_{\rm lead}$ uses only the uniform free moments and the
fixed derivative and reset conventions.  It does not contain
a constant that grows with $L$.  Choose the seven canonical
caps successively using Lemma~\ref{rg:canonical-contraction}.
Choose fixed triangular weights for their discrepancy norm.

Let $D$ exceed every fixed polynomial degree in $M$, support
load, Gaussian score variables, and normalization entropy
used above.  Increase it for the finitely many remaining
regroupings and choose $P_0\ge2D+2$.  Finally increase $H$ until,
uniformly for $H_j\ge H$,
\begin{align}
 C_Le^{-c_LM}\operatorname{poly}(M)&\ll1,&
 C_LtM^D&\ll1,\nonumber\\
 C_LM^D/p^2&\ll1,&
 C_LM^D(T^2/t+e^{-c_LM}/t+t)&\ll1,\nonumber\\
 C_LM^D(1+\log g^{-1})&<c_Lp^2/4,&
 c_Lp^2/M^D&\gg p^{3/10}+D\log g^{-1}+D\log M,\nonumber\\
 C_Lg^{.04}\operatorname{poly}(M,p)&\ll M^{-2},&
 C_Lg^{4.6}/\delta_j&\ll1 .
 \label{rg:finite-choices}
\end{align}
Each condition follows either from $P_0\ge2D+2$ or from a
positive power of $H_j$ dominating every fixed power of its
logarithm.  For instance $T^2/t=H_j^{-.48}/p$, and
$g^{4.6}/\delta_j=O(H_j^{-.25})$.  The choices are therefore
noncircular.

Equation~\eqref{rg:closed-error} now gives the renewed error cap.
Every output-seeded component retains a compulsory suppression.
The reserve, the shared Gaussian tail estimate and
\eqref{rg:finite-choices} give, after all its decorations and
regroupings, the stronger covering bound
$e^{-p^{3/10}}=o(w_{j-1})$.  This pays the fixed hard-derivative
losses and the relative input error and gas discrepancy norms.
Symmetrizing the assignments by the finite spatial group
preserves these bounds and the density.

For completeness, divide the error discrepancy in
\eqref{rg:raw-error} by $\delta_{j-1}$.
The remainder factor is small because
$g^{4.5}/\delta_j=O(g^{.4})$.  The gas has just been bounded with
strictly more suppression than its required relative cap.
The canonical map is triangular with contracting diagonal.
Choose the later discrepancy weights sufficiently small, and
then enlarge $H$ for the remaining cross terms.  This gives a
common $q<1$.  Dependence on the precise coupling costs
$\operatorname{poly}(M,p)|b_2-b_1|/H_j$, which is at most
$C_LH_j^{-1/2}|b_2-b_1|$ after increasing $H$.
Differences of the canonical kicks cost $C_Lu$; differentiating
$\kappa/b$ costs $C_L|b_2-b_1|/H_j^2$; the extracted raw Hessian
has the bounds just proved.  Weakening the positive decay
exponent if necessary gives
\eqref{rg:shape-comparison}--\eqref{rg:coupling-comparison}.
This finishes the proof of Theorem~\ref{thm:rg-closure}.
